\chapter{Quantum--tropical fluid realizations: operators, control and substrates}
\label{ch:quantum-tropical-realizations}

\reviewstatus{This chapter develops an explicitly proposed realization of
ToricGT-style graph models, together with finite mathematical results and an
exact two-qubit feedback example. Its operator, tropical and control layers
have separate contracts. The physical literature supplies demonstrated
building blocks, not an experimental realization of the combined architecture.
Generative substrate design and the mappings to natural systems below are
research proposals. The new finite arguments are written proofs; only the
weighted-gain components identified in the formal type map have Lean coverage.
No new fluid-trained model, fabricated substrate or quantum speedup is reported.}

\section{What a realization must specify}

Here a \emph{quantum--tropical fluid realization} means a model with three
explicit ingredients: quantum states and admissible operations on them;
graph-based routing or a logarithmic asymptotic description; and dynamics
whose relation to a fluid or other physical substrate is quantitatively
specified. It is a proposed interface, rather than a claim that these three
descriptions are interchangeable. A useful realization must identify an
encoding, an evolution, a readout and the approximation error at each
interface.

The reviewed ToricGT implementation contributes graph tokens, endpoint
incidence, finite Weyl operators, a mixture of ordinary and max-plus
attention, and expert routing. The program and Tropical Quivers manuscripts
add proposed algebraic and dynamical constraints. Their exact revisions and
the corrections to selected statements are recorded in
Chapter~\ref{ch:toric-transport-audit}. We retain that pinned meaning of
``ToricGT-style''; the name alone imposes neither a hydrodynamic limit nor a
physical quantum representation.\cite{qtToricGTCode,qtToricGTProgram,ToricAuditQuivers2026}

A concrete specification can be written
\[
 (Q,\mathcal R,\mathcal A,\mathcal E,\Phi,\mathcal M,\pi,\mathcal P).
\]
Here $Q$ is a finite quiver, $\mathcal R$ its intended relations,
$\mathcal A$ the operator algebra, $\mathcal E$ an encoding,
$\Phi$ an admissible evolution, $\mathcal M$ a measurement instrument,
$\pi$ a classical or coherent control law, and $\mathcal P$ a physical
model with units, boundary conditions and parameter uncertainty. A compiler
must say which of these structures it preserves. The fluid discretization
and its continuum consistency belong to $\mathcal P$; the predictor's
training objective belongs to a further statistical specification.

For example, a gentle-algebra relation declaring a composition of two
arrows zero is meaningful for a representation by linear maps on vector
spaces. A composition of trace-preserving channels on a nonzero state space
cannot be the zero map: it still preserves trace. Thus a forbidden path may
be enforced in a classical routing graph, or represented by trace-decreasing
branches of an instrument, but cannot simply be imposed as a zero relation
between unconditional trace-preserving channels. Relations must be checked
in the category in which the proposed hardware actually operates.

\section{A finite operator algebra with conservative quantum transport}

\subsection{Classical position and quantum internal state}

Assign a finite-dimensional Hilbert space $\HH_i$ to each vertex and set
\[
 \mathcal A_Q=\bigoplus_{i\in V}\mathcal B(\HH_i).
\]
A state is a tuple $\rho=(\rho_i)$ with $\rho_i\succeq0$ and
$\sum_i\Tr\rho_i=1$. The mass at vertex $i$ is
$m_i=\Tr\rho_i$; when $m_i>0$, $\rho_i/m_i$ is its conditional internal
state. The direct sum describes a classical location register with quantum
internal degrees of freedom. Coherences between locations instead require
the full algebra on $\bigoplus_i\HH_i$, including off-diagonal blocks.
This choice changes the state space and therefore changes the physics.
Neither finite description is automatically a many-body quantum fluid:
a tensor-product or Fock-space model has additional particle and correlation
structure. An encoding into finitely many blocks must state which
observables it retains and bound truncation or leakage effects.

Let $P_{ij}\ge0$, $\sum_iP_{ij}=1$, and let
$\Phi_{ij}:\mathcal B(\HH_j)\to\mathcal B(\HH_i)$ be completely
positive trace-preserving maps for each allowed transition. Define
\begin{equation}\label{eq:qtf-transport}
 (\mathcal T\rho)_i=\sum_jP_{ij}\Phi_{ij}(\rho_j).
\end{equation}
This finite construction is related to open quantum walks, whose
operator-valued transition rules combine quantum internal states with
classical position records.\cite{qtsAttal2012}

\begin{proposition}[Conservative transport on a direct-sum algebra]
\label{prop:qtf-channel-transport}
Under the preceding assumptions, $\mathcal T$ is completely positive and
trace preserving. Its masses satisfy $m'=Pm$. For Hermitian inputs
$X=(X_j)$,
\[
 \sum_i\|(\mathcal TX)_i\|_1\le\sum_j\|X_j\|_1.
\]
In particular the total trace distance between two states cannot increase.
\end{proposition}
\begin{proof}
Write $\Phi_{ij}(X)=\sum_\alpha K_{ij\alpha}XK_{ij\alpha}^*$ with
$\sum_\alpha K_{ij\alpha}^*K_{ij\alpha}=I_{\HH_j}$. Regard each
$\sqrt{P_{ij}}K_{ij\alpha}$ as a block operator supported from $j$ to
$i$. Their Kraus sum defines a completely positive map, including after
tensoring with any finite ancillary matrix algebra. On input block $j$,
the sum of their adjoint products is
$\sum_iP_{ij}I_{\HH_j}=I_{\HH_j}$, which proves trace preservation.
Taking the trace of each output block gives $m'=Pm$.

For a Hermitian $X_j$, write its positive and negative parts as
$X_j=X_j^+-X_j^-$. Positivity and trace preservation imply
\[
 \|\Phi_{ij}(X_j)\|_1
 \le\Tr\Phi_{ij}(X_j^+)+\Tr\Phi_{ij}(X_j^-)
 =\|X_j\|_1.
\]
The triangle inequality and column sums of $P$ then give
\[
 \sum_i\Bigl\|\sum_jP_{ij}\Phi_{ij}(X_j)\Bigr\|_1
 \le\sum_{i,j}P_{ij}\|X_j\|_1
 =\sum_j\|X_j\|_1.
\]
\end{proof}

Complete positivity is an algebraic constraint on a map and its extensions
to entangled ancillary systems. A numerical test on a few positive input
matrices does not establish it. A Kraus representation with a checked
normalization, or an exactly positive Choi matrix with the correct partial
trace, gives a finite certificate. The Heisenberg dual $\mathcal T^*$ acts
linearly on observables and is unital. These statements concern the same
channel in two dual descriptions.

\subsection{The extra requirements imposed by a fluid interpretation}

A local transition matrix depending on a prescribed velocity field can
transport the masses $m_i$. Adding internal channels can represent
collisions, relaxation or a controlled dissipative feature map. This is a
finite transport model. To obtain a specified fluid equation one must also
identify physical density, momentum, energy, length and time scales; prove
the required conservation laws; specify constitutive assumptions; and
control the discretization and kinetic limits. A positive mass-conserving
map alone has no specified pressure law, viscosity or incompressibility.
The reduction contracts in the preceding chapters state the corresponding
closure and convergence obligations.

For a continuous-time finite quantum model, a standard admissible generator
has the form
\[
 \dot\rho=-i[H,\rho]+\sum_\alpha
 \left(L_\alpha\rho L_\alpha^*
 -\tfrac12\{L_\alpha^*L_\alpha,\rho\}\right),
 \qquad H=H^*.
\]
With finite matrices and time-independent coefficients, this generates a
completely positive trace-preserving semigroup.\cite{qtsLindblad1976}
Local choices of the $L_\alpha$ encode dissipative graph transitions;
off-diagonal Hamiltonian couplings can support coherent transport. A forward
Euler step need not itself be completely positive, even when the continuous
generator is admissible. A numerical realization must retain that
distinction or use a structure-preserving step.

Quantum fluids of light provide a particularly relevant physical family:
interactions, coherent driving and loss lead to rich collective dynamics,
with different descriptions appropriate to mean-field and correlated
regimes. Their equations supply a possible physical model for coupled
tokens; a generic optical neural-network nonlinearity does not automatically
reproduce the correlated quantum regime.\cite{qtsCarusottoCiuti2013,qtsAmo2009}

\section{Toric features, tropical routing and the information they discard}

On $\C^d$, define $X|k\rangle=|k+1\bmod d\rangle$ and
$Z|k\rangle=\omega^k|k\rangle$, where $\omega=e^{2\pi i/d}$. Then
$ZX=\omega XZ$. Products $Z^aX^b$ give finite Weyl features, while
expectations $\Tr(O\rho)$ give real features for Hermitian $O$.
A unitary change of basis transports both $O$ and $\rho$ and leaves their
pairing unchanged. A learned interface may therefore use these finite
operator features while retaining a classical graph controller. Additional
claims about a toric variety, a torus action or a BGG complex require their
own explicitly preserved structure.

The tropical part can enter as a routing optimizer or as a logarithmic
envelope. For $\tau>0$ and $m\ge1$ real scores $s_j$, define
\[
 L_\tau(s)=\tau\log\sum_{j=1}^m e^{s_j/\tau},\qquad
 p_j^\tau=\frac{e^{s_j/\tau}}{\sum_k e^{s_k/\tau}}.
\]

\begin{proposition}[Controlled tropical limit of a positive sum]
\label{prop:qtf-tropical-limit}
For $M=\max_js_j$,
\[
 M\le L_\tau(s)\le M+\tau\log m.
\]
If the maximizer $j_*$ is unique and
$s_{j_*}-s_j\ge\gamma>0$ for all $j\ne j_*$, then
\[
 p_{j_*}^\tau\ge
 \frac{1}{1+(m-1)e^{-\gamma/\tau}}.
\]
\end{proposition}
\begin{proof}
Factor $e^{M/\tau}$ from the sum. The remaining sum is between one and
$m$, giving the first bounds after taking its logarithm. Divide the
normalizing sum for $p_{j_*}^\tau$ by $e^{s_{j_*}/\tau}$.
Each other term is at most $e^{-\gamma/\tau}$, which gives the second
bound.
\end{proof}

This statement concerns a positive exponential sum. With several maximizers
the limiting probability is supported on that set; there is no uniquely
selected route without an extra rule. The reviewed tropical-attention code
uses scores that depend on dot products and applies a max-plus aggregation.
Even when the aggregation is piecewise affine in its supplied scores,
the complete map need not be piecewise affine in all original inputs.

For a product of nonnegative gain matrices, logarithms turn multiplication
along one path into addition. Taking the largest path gives a max-plus
envelope, with $\log0=-\infty$ and zero-weight paths absent. Ordinary
matrix multiplication also sums over paths; its norm
requires the path-multiplicity factor proved in
Chapter~\ref{ch:toric-transport-audit}. Signed or complex amplitudes add
phase cancellation as another issue. A tropical magnitude envelope can
upper-bound selected effects once these factors are supplied, but does not
retain the interference pattern or determine entanglement.

Min-plus arithmetic also arises directly in control. A deterministic
finite-horizon route problem has the Bellman recursion
\[
 V_n(i)=\min_{a\in\mathcal U(i)}
 \{c_n(i,a)+V_{n+1}(f(i,a))\}.
\]
This is a precise optimization layer that can select graph operations.
A stochastic, risk-sensitive or entropy-regularized controller has a
different recursion, whose transition probabilities and normalization must
be retained. A zero-temperature approximation can reduce exploration but
also creates switching boundaries. The common gain certificate below is
useful only if it covers every allowed branch and the effects of switching.

\section{Nonlinear dynamics on an operator state space}

\subsection{Where the nonlinearity lives}

An operator algebra is a linear space with a product and an adjoint. There
are several mathematically distinct ways to obtain nonlinear dynamics
associated with it. First, a classical controller $z$ can select a channel
or a generator:
\[
 \rho_{n+1}=\Phi_{\theta(z_n)}(\rho_n),\qquad
 z_{n+1}=f(z_n,y_n),\qquad
 y_n=\mathcal M_n(\rho_n).
\]
For a fixed classical setting, $\Phi_{\theta}$ is linear on density
operators. The coupled controlled system can be nonlinear. If $y_n$
denotes an ideal expectation rather than an actual measurement record,
the displayed recursion is an ensemble-level model and needs a readout
approximation before it describes hardware.
For actual measurement records, the conditional measurement backaction
must be included in the state update, or readout must use separately
prepared copies. It cannot be omitted while retaining a claim about the
same undisturbed quantum system.

Second, a measurement instrument has unnormalized branches
$\rho\mapsto K_y\rho K_y^*$. Each is linear and completely positive.
Conditioning on a nonzero-probability outcome divides by
$\Tr(K_y\rho K_y^*)$ and produces a nonlinear normalized update. A
physically complete account retains the probabilities and the classical
outcome register. Replacing that instrument with a deterministic normalized
map on an unknown state changes the operation.

Third, effective mean-field equations may contain a state-dependent
Hamiltonian, for example $\dot\rho=-i[H(\rho),\rho]$, or a cubic
order-parameter equation. Such equations approximate selected observables
or many-body regimes under stated assumptions. Their nonlinearity does not
by itself supply a fundamental nonlinear quantum channel. Driven polariton
and atomic condensate models make this distinction experimentally relevant.
\cite{qtsCarusottoCiuti2013,qtsBarontini2013}

These three cases give a practical interpretation of a nonlinear operator
perspective: learn and control dynamics \emph{on} a physically admissible
state space while retaining the linear completely positive operations
appropriate to each fixed experimental setting. Positivity, trace, symmetry,
causality and accessible controls become architectural constraints.

\subsection{Koopman lifting and its closure error}

For a deterministic state-space map $F:\mathcal X\to\mathcal X$, the
Koopman operator is $(Uf)(x)=f(F(x))$. It is linear in the observable $f$
even when $F$ is nonlinear. On an algebra of ordinary scalar observables it
also preserves products. For a stochastic transition, conditional
expectation instead gives a positive unital operator, generally without
product preservation. Neither construction is automatically the
Heisenberg dual of a channel on the noncommutative algebra $\mathcal A_Q$:
their domains and their notions of positivity must first be matched.

A finite dictionary $\psi(x)$ may satisfy
\[
 \psi(F(x))=K\psi(x)+r(x).
\]
An exact finite representation requires $r=0$ on the claimed domain.
Learning $K$ from trajectories ordinarily leaves a residual, and a small
training residual is not a uniform domain bound. Finite invariant
subspaces that also reconstruct the state impose substantial restrictions,
as the primary Koopman-control analysis explains.\cite{qtsBrunton2016}
The residual belongs in an error recurrence such as
Theorem~\ref{thm:weighted-gain-propagation}. This connects operator learning
to a checked control budget without declaring a finite closure exact.

\section{Control, observability and robustness}

\subsection{A common gain for coupled classical and quantum blocks}

Let $d_\rho$ measure quantum-state error, for example trace distance, and
let $d_z$ measure controller error. Suppose the actual closed-loop map
satisfies on an invariant domain
\[
 \begin{pmatrix}d_\rho'\\d_z'\end{pmatrix}
 \le
 G\begin{pmatrix}d_\rho\\d_z\end{pmatrix}
 +\begin{pmatrix}\eta_\rho\\\eta_z\end{pmatrix},\qquad
 G=\begin{pmatrix}a&b\\c&d\end{pmatrix}\ge0.
\]
For $q>\max(a,d)$ and $b,c>0$, the certificate
$Gw\le qw$ with $w_\rho,w_z>0$ is equivalent to
\[
 \frac{b}{q-a}\le\frac{w_\rho}{w_z}
 \le\frac{q-d}{c}.
\]
Indeed the first row rearranges to the left inequality and the second to
the right. Feasibility is therefore equivalent to
$bc\le(q-a)(q-d)$ at that rate. If $a,d<1$ and
$bc<(1-a)(1-d)$, continuity supplies a feasible $q<1$ sufficiently close
to one. Cases with zero coupling can be checked directly using $Gw\le qw$.
These formulas expose the cost of feedback: individually stable components
can still have an excessive loop gain.

Once the gains are justified, a common $w$ and $q$ give the full finite-time
budget of Theorem~\ref{thm:weighted-gain-propagation}. An optimizer may
propose $w$; the independent rational checker verifies the proposed
inequalities. Separate certificates can cover model discrepancy, channel
approximation, integration error and readout error as explicit residuals.
For a uniform weighted perturbation $\delta$ and $q<1$, the eventual
bound is $\delta/(1-q)$. A smaller $q$ improves this robustness bound,
while rapid contraction also suppresses memory of past inputs. Reservoir
design therefore needs a task-dependent tradeoff between memory, separation
of inputs and noise sensitivity.\cite{qtsFujiiNakajima2017}

\subsection{Accessible controls and finite readout}

For a closed finite quantum system, accessible Hamiltonians generate a Lie
algebra that constrains the reachable unitaries. Full unitary control still
preserves the eigenvalues of a density matrix. Changing its spectrum
requires measurement, dissipation or an enlarged system and subsequent
discarding. Thus the ability to fit a target matrix numerically does not
prove that a specific device can implement it.\cite{qtsAlbertiniDAlessandro2003}
For an open device, the allowed jump operators, rates, locality, noise and
bandwidth similarly constrain attainable channels.

Observability has a corresponding cost. One cannot obtain every expectation
of an unknown state from one undisturbed specimen. A hardware proposal must
specify repeated preparations, ensemble access, weak measurement or an
instrument with a state estimator. A statistical error guarantee over a
finite number of measurements can be converted into a deterministic
residual bound on the event that those estimates are accurate. The event's
confidence, the number of shots and the finite time horizon remain part of
the result; an unlimited-horizon guarantee does not follow automatically.

For example, suppose $M$ measurement settings are fixed in advance and
each uses $N$ independent shots with outcomes in $[-1,1]$. Hoeffding's
inequality and a union bound give simultaneous error at most $\varepsilon$
with probability at least $1-\delta$ when
\[
 N\ge\frac{2}{\varepsilon^2}\log\frac{2M}{\delta}.
\]
For one setting, the two-sided failure probability is at most
$2e^{-N\varepsilon^2/2}$; adding these $M$ bounds proves the assertion.
The confidence concerns sampling error for those fixed settings, not
calibration bias or arbitrary temporally correlated reservoir outputs.
Adaptive feedback requires a corresponding conditional sampling argument
and a specified finite horizon.\cite{qtsHoeffding1963}

Model predictive control can optimize a finite horizon subject to channel
admissibility, actuator limits and a certified invariant set. For a
classical energy model
\[
 \dot x=(J(x)-R(x))\nabla E(x)+B(x)u,\qquad
 J^T=-J,\quad R\succeq0,
\]
the chain rule gives
$\dot E=-\nabla E^TR\nabla E+y^Tu$ with
$y=B^T\nabla E$. This is a useful energy-balance constraint for a
fluid, elastic or electrical surrogate. A quantum dissipative model needs
its own thermodynamic assumptions: an arbitrary driven channel need not
decrease a chosen energy or increase entropy. The proposed controller
must preserve the balance appropriate to its physical model.

\section{An exact nonlinear quantum-state feedback example}

The following small construction instantiates the control contract without
assuming a gain bound from sampled derivatives. It uses a tuple of two
qubit density matrices; it does not encode entanglement between the two
blocks. Let $Z=\diag(1,-1)$,
\[
 z_j=\Tr(Z\rho_j),\quad
 B=\begin{pmatrix}0&2\\1/8&0\end{pmatrix},\quad
 b=\begin{pmatrix}1/4\\0\end{pmatrix},\quad
 u_i=\operatorname{clip}((Bz+b)_i,-1,1).
\]
Set $a=1/4$, $\tau(u)=(I+uZ)/2$, and define
\begin{equation}\label{eq:qtf-feedback}
 F_i(\rho)=a\rho_i+(1-a)\tau(u_i).
\end{equation}
The clipped controller is nonlinear but globally Lipschitz. Its expectations
are assumed known exactly for the ideal model. A device implementation
would estimate them through the readout process just described.

\begin{proposition}[A physical channel family and a contracting feedback law]
\label{prop:qtf-feedback-example}
For every fixed $u\in[-1,1]$,
\[
 \Phi_u(X)=aX+(1-a)\tau(u)\Tr X
\]
is completely positive and trace preserving. The ideal feedback map $F$
preserves pairs of density matrices and satisfies, with
$D_i=\tfrac12\|\rho_i-\sigma_i\|_1$,
\[
 D_i'\le\sum_jG_{ij}D_j,\qquad
 G=\begin{pmatrix}1/4&3/2\\3/32&1/4\end{pmatrix},\qquad
 G\begin{pmatrix}4\\1\end{pmatrix}
 =\tfrac58\begin{pmatrix}4\\1\end{pmatrix}.
\]
It therefore contracts the weighted maximum trace distance at rate $5/8$.
Its fixed point has Bloch vectors $(0,0,1/3)$ and $(0,0,1/24)$.
The feedback map itself is not affine in the pair of input states.
\end{proposition}
\begin{proof}
Since $|u|\le1$, $\tau(u)$ is a density matrix. In input--output order,
the unnormalized Choi matrix is
\[
 J(\Phi_u)=a|\Omega\rangle\langle\Omega|
 +(1-a)I\otimes\tau(u),\qquad
 |\Omega\rangle=|00\rangle+|11\rangle.
\]
Both terms are positive, and tracing over the output gives
$aI+(1-a)I=I$. Equivalently the replacement map has Kraus operators
$\sqrt{t_k}|k\rangle\langle l|$, where $t_k$ are the diagonal entries
of $\tau(u)$ and $k,l\in\{0,1\}$. Scaling these replacement operators
by $\sqrt{1-a}$ and adjoining $\sqrt a I$ gives the same channel.
Convexity then proves state-space
invariance for every chosen control.

Measurement duality gives $|\Delta z_j|\le2D_j$. Clipping is
one-Lipschitz and
$\tfrac12\|\tau(u)-\tau(v)\|_1=|u-v|/2$. Thus
\[
 D_i'\le aD_i+\frac{1-a}{2}|\Delta u_i|
 \le aD_i+(1-a)\sum_j|B_{ij}|D_j.
\]
Substitution gives $G$, and direct multiplication gives its certificate.
The state space is complete in this finite-dimensional metric, so the
strict contraction has at most one fixed point and converges to it.
At a fixed point the transverse Bloch coordinates vanish; the displayed
$z$ values solve $z=Bz+b$ and lie inside the clipping interval. They
therefore give existence as well as uniqueness.

For nonaffinity, hold the first input at $z_1=0$ and choose second inputs
with $z_2=0$ and $z_2=1$. The first output $z$ values are $3/16$ and
$3/4$. Their average is $15/32$. At the midpoint input $z_2=1/2$,
the clipped control is one and the first output is $3/4$. The difference
is $9/32$, so a single affine channel cannot equal this ideal feedback
map on all inputs.
\end{proof}

The fixed-control channels are physically admissible. The ideal feedback
recursion describes what happens when a controller has the required
expectations; it is not an instruction for obtaining them without readout
cost. The associated exact implementation uses rational Bloch coordinates
and rational Choi matrices. It retains a coherent initial state, tests an
actually saturated step, and checks trace-distance bounds by squaring
nonnegative quantities, avoiding floating-point square-root tolerances.
These replacement channels are generally nonunital:
$\Phi_u(I)=I+(1-a)uZ$. The entropy-monotonicity statement for the earlier
unital heat channel consequently does not transfer to this example.

There is also an exact tropical comparison. Weighted normalization gives
\[
 K=\diag(w)^{-1}G\diag(w)
 =\begin{pmatrix}1/4&3/8\\3/8&1/4\end{pmatrix}.
\]
Every row of $K^n$ sums to $(5/8)^n$, whereas the largest weight of one
length-$n$ path is $(3/8)^n$. The latter is attained by alternating the
vertices. The difference comes from adding the other paths, exactly as in
the path-multiplicity analysis. The proposed rate $3/5$ fails the supplied
weighted certificate and is retained as negative replay evidence.

Finally, if each measured expectation has deterministic error at most
$\varepsilon$, the extra weighted one-step discrepancy is at most
$3\varepsilon/16$: the two component bounds after division by $w_i$
are $3\varepsilon/16$ and $3\varepsilon/64$. Therefore the same
residual theorem gives the asymptotic envelope
\[
 \frac{3\varepsilon/16}{1-5/8}=\frac{\varepsilon}{2}.
\]
For $\varepsilon=1/100$, this is $1/200$. This is conditional error
control under the stated readout bound. It is not an experimentally
measured noise floor or a substitute for a finite-shot confidence analysis.

\section{Candidate engineered and natural substrates}

A substrate is relevant when its controlled degrees of freedom implement
the specified encoding, operation and readout with bounded error. Its
microscopic quantum origin alone does not establish useful coherent quantum
processing. The following examples identify both the usable mechanism and
the experiment needed to connect it to the proposed architecture.

\begin{longtable}{p{0.20\textwidth}p{0.34\textwidth}p{0.36\textwidth}}
\toprule
Substrate & Candidate role & Evidence and realization test\\\midrule
\endhead
Programmable photonic networks &
Interferometers implement coherent mode mixing; phase shifters and
detectors can supply a graph operator and readout. A separate controller
can update routing. &
The cited programmable-transport experiment used continuous-wave laser
illumination and intensity readout. A few-photon channel would need its
own nonclassical-state specification and validation, alongside loss and
control-drift measurements.
\cite{qtsHarris2017}\\[6pt]
Exciton--polariton fluids &
Coupled cavities or spatial modes supply driven, lossy, interacting
fields. They are candidates for nonlinear reservoirs and graph-coupled
fluid features. &
Quantum-fluid and superfluid phenomena motivate the mechanism. A
realization must identify its mean-field or correlated regime, its
input/output map, and the controller's error bounds.
\cite{qtsCarusottoCiuti2013,qtsAmo2009}\\[6pt]
Ultracold atomic gases &
Trapped modes or lattice sites provide quantum matter with controllable
couplings and, in selected settings, engineered local dissipation. &
Dissipative many-body dynamics have been studied experimentally. A
ToricGT mapping would additionally need accessible graph couplings,
repeated state preparation and an observable task with a calibrated
readout budget.\cite{qtsBarontini2013}\\[6pt]
Superconducting cavity circuits &
Microwave modes, nonlinear elements and engineered loss can implement
bosonic channels, memory and controlled dissipation. &
Engineered dissipation has confined dynamics to a quantum-state manifold.
The proposed graph realization must quantify truncation, leakage,
decoherence and the attainable channel family.\cite{qtsLeghtas2015}\\[6pt]
Magnonic and magnetic networks &
Spin-wave propagation and nonlinear responses can provide a compact
dynamical reservoir with graph-like couplings. &
Magnetic reservoir experiments support this computing mechanism. The
usual wave-regime demonstration establishes a physical reservoir;
quantum-state processing would need its own preparation and coherence
evidence.\cite{qtsNamiki2023}\\[6pt]
Soft bodies and fluid--structure systems &
Elastic motion, damping and hydrodynamic coupling can distribute
nonlinear memory across a body, with sensors and actuation supplying
the interface. &
Soft-body reservoir experiments demonstrate task-relevant dynamics.
Natural morphologies motivate graph constraints, but every generated
body needs a measured transfer map and robustness test.
\cite{qtsNakajima2015}\\[6pt]
Photosynthetic molecular networks &
Excitonic couplings, vibrations and environmental relaxation offer a
natural example of open-system energy transport. &
The functional role and character of observed coherence require care.
Electronic and vibrational contributions, disorder and open-system
conditions must be distinguished before proposing an information
processor.\cite{qtsDuan2017,qtsHarushDubi2021}\\\bottomrule
\end{longtable}

Natural vascular, porous and transport networks provide another source of
geometries and conservation constraints for a classical graph-fluid
realization. Their branching and feedback suggest a design grammar;
they do not supply a trained ToricGT model by identification alone.
Likewise, a crystal or a biological assembly discovered in nature can
inspire a coupling pattern without providing arbitrary programmable
control. The useful question is which functions the accessible dynamics
can implement and which invariants survive the readout.

There is no universal ordering of these substrates. Interferometric
systems favor controlled linear mixing; strongly interacting quantum
platforms favor selected many-body phenomena; dissipative mechanical and
magnetic reservoirs may provide economical nonlinear memory. A comparison
needs a common task, input bandwidth, output accuracy, reset protocol,
energy accounting and training budget. Device operation, control and
measurement must all enter the cost.

\section{De novo substrate design with generative models}

\subsection{Generate a physical candidate for a specified operator}

A generative model can propose a material or geometry conditioned on a
target response. Inorganic crystal generation and generative photonic
inverse design provide concrete examples of this approach. Their outputs
are candidates for a physical validation pipeline; the generative likelihood
is not an operator-realization certificate.\cite{qtsZeni2025,qtsJiangFan2019}

Let $s$ describe composition, geometry, connectivity and fabrication
parameters, and let $u$ describe accessible controls. A forward model
produces an effective evolution $\mathcal F_{s,u}$ and readout
$\mathcal M_s$. One can seek
\[
 \min_{s,u,w,q}\quad
 L_{\rm task}(s,u)+\lambda_E E(s,u)+\lambda_C C(s)
 +\lambda_R R_{\rm uncertainty}(s,u)
\]
subject to physical admissibility, an observable target interface, and
$G(s,u)w\le qw$ with $w>0$, $q<1$ when contraction is appropriate.
The objective terms need measured definitions: prediction error,
energy per usable output, fabrication or control cost, and a robustness
penalty. A generative prior can constrain plausible $s$ without replacing
the physical constraints.

The specification should begin with a target channel, scattering map,
response kernel or controlled fluid update, rather than with a visual
shape. For an optical device, relevant constraints include a realizable
material response, loss, bandwidth and fabrication tolerances. For a
magnetic, elastic or fluid device, stability, conservation laws, boundary
conditions and sensor coupling matter. For a proposed quantum device,
state preparation, available interactions and admissible measurement are
part of the target. Complete positivity constrains a reduced channel;
additional locality and controllability constraints determine whether the
chosen substrate can realize it.

\subsection{A falsifiable co-design loop}

The following sequence separates design exploration from validation.
\begin{enumerate}
\item Specify the graph task, allowed algebraic relations, encoding,
observables, error tolerance and resource budget. Decide whether the
target requires coherent quantum states or only a wave/dynamical response.
\item Generate candidate structures conditioned on those requirements.
Apply exact geometry, connectivity and composition constraints where
available, and retain the generator version and random seed.
\item Evaluate candidates with the appropriate physical solver and
uncertainty model. Reject violations of conservation, positivity,
stability or attainable controls; a good surrogate score cannot override
these checks.
\item Identify a reduced input--state--output model on a stated domain.
Measure its residual on held-out conditions and obtain a justified gain
envelope. If only sampled derivatives are available, record that empirical
scope rather than promoting it to a uniform bound.
\item Co-design the controller and positive weights. Replay finite
rational certificates independently, and propagate solver, readout and
fabrication errors through the common-gain budget.
\item Evaluate prospective fabrication variation, then fabricate or
assemble a selected candidate only under an actual experimental program.
Calibrate it, test closed-loop behavior and compare with the predicted
operator. A simulation-only result remains labeled as such.
\item Compare against conventional numerical solvers and simpler
physical or classical learning baselines on the same task. Include
state preparation, measurement, reset, calibration and controller costs.
\end{enumerate}

MatterGen establishes that conditional generative modeling can propose
inorganic crystals with target properties, including an experimentally
tested candidate. It does not supply the mapping from a desired quantum
channel to a fabricable quantum device.\cite{qtsZeni2025}
Generative photonic inverse design directly targets an optical response,
but its particular demonstrated response still needs an additional mapping
to this graph-fluid architecture.\cite{qtsJiangFan2019}
Those missing mappings are useful research questions: they specify what a
new project must actually demonstrate.

Natural systems can enter this loop as a prior on topology, material motifs
or dynamics. A learned grammar might favor sparse transport networks,
hierarchical elastic couplings or molecular geometries. The design must
then be evaluated with the same operator, control and readout contracts
as a wholly synthetic candidate. This makes biomimetic inspiration
testable without treating an analogy as a realization.

\section{Connections to the wider research program}

The construction offers a practical role for the project's order and
surface-algebra language: encode legal compositions and symmetries, select
observable features, and demand that a realization preserve specified
relations. Quantum information supplies admissible state transformations
and operational distances. Tropical mathematics supplies routing and
logarithmic bounds. Control theory connects these finite operations to
robust dynamical behavior. Each connection has an explicit map and an
error obligation.

For Hilbert--P\'olya research, a controllable finite operator family can
test spectral statistics or evaluate a proposed discretization. A finite
fit to zeros does not establish the required self-adjoint operator,
spectral identification or analytic limiting theorem. For Yang--Mills,
an engineered finite dissipative or Hamiltonian gap is a property of that
system and has no automatic implication for the continuum mass-gap
problem. For cryptography, a generated physical response may be a useful
device primitive, but a security claim additionally requires a threat
model, entropy or hardness assumptions, extraction and implementation
analysis. The existing cryptographic implementations and their security
scope remain separate evidence.

For Navier--Stokes, the useful next experiments are bounded: preservation
of discrete invariants, closure-error prediction, controllable manufactured
solutions and comparisons with conventional solvers at stated resolution.
The externally released forced singularity theorem retains its own
continuum hypotheses and independent-build status. The present finite
control construction neither supplies a stronger continuum theorem nor
removes the unresolved independent upstream build.

The immediate implementation milestone is the exact feedback example and
its independent gain replay. The next research milestone is a trained
graph-fluid surrogate with a documented invariant domain and held-out
error budget. A physical milestone requires one specified substrate,
calibrated encoding and readout, measured operator agreement and a fair
resource comparison.
